% ============================================================
%  style/griffe.tex — \griffe: Griffbilder unter/neben dem Songtext platzieren
%  (wird von style.tex eingebunden)
% ============================================================
% ------------------------------------------------------------
% \griffe[Griffe pro Zeile]{F, A\#, Bb/D}
% Griffbilder (aus chords.tex, griffe-zusatz.tex oder \defgriff), grau.
% Platzierung MANUELL, wie bei \schlagmuster:
%  - Hinten an eine Zeile gehängt -> rechts am Textrand auf Höhe dieser
%    Zeile, über den Text gelegt (nach unten wachsend), z.B.
%      \slbl{Chorus:}\griffe{F, C/E}
%      \cl{Am   F}\griffe[2]{Fmaj7, Am7, Dsus2}   (2 Griffe pro Zeile)
%  - Allein in einer eigenen Zeile -> wie bei \schlagmuster eine Leerzeile,
%    die Griffe liegen rechts am Rand auf Höhe dieser Leerzeile.
%  - Die Griffe selbst schieben nie Text weg (werden nur darübergelegt).
%  - Griffe pro Zeile optional (Standard: alle nebeneinander, so viele wie
%    in die Breite passen).
% Unbekannte Namen -> leeres Griffbrett + Warnung.
% ------------------------------------------------------------
\makeatletter

% --- Positionen der Songzeilen/Schlagmuster in die .aux schreiben ---
% Schlüssel: Song-Nr. + Seite innerhalb des Songs (stabil, auch wenn sich
% die Seitenzahlen zwischen zwei Durchläufen verschieben).
\newcounter{griffsong}
\newcommand{\griffsongkey}{0}
\newcommand{\griffsongstart}{0}
\newcommand{\griffsongbeginn}{%
  \stepcounter{griffsong}%
  % fester Schlüssel pro Song (unabhängig von der Buchreihenfolge)
  \ifx\aktuellesongid\empty
    \xdef\griffsongkey{n\arabic{griffsong}}%
  \else
    \xdef\griffsongkey{\aktuellesongid}%
  \fi
  \xdef\griffsongstart{\the\c@page}%
}
\newcommand{\griffleftodd}{0}\newcommand{\griffleftevn}{0}\newcommand{\griffbottom}{0}
\AtBeginDocument{%
  \xdef\griffleftodd{\number\dimexpr 1in+\hoffset+\oddsidemargin\relax}%
  \xdef\griffleftevn{\number\dimexpr 1in+\hoffset+\evensidemargin\relax}%
  \xdef\griffbottom{\number\dimexpr\paperheight-1in-\voffset-\topmargin
                     -\headheight-\headsep-\textheight\relax}%
}
\newcommand{\griffzeileoben}{524288}   % 8pt über der Grundlinie (in sp)
\newcommand{\griffzeileunten}{196608}  % 3pt unter der Grundlinie (in sp)

% #1 rechtes Ende (x relativ zum linken Textrand, sp)
% #2 / #3 Ausdehnung über / unter der aktuellen Grundlinie (sp)
\newcommand{\griffschreiben}[3]{%
  \pdfsavepos
  \edef\grifftmpwrite{\write\@auxout{\string\griffhindernis{\griffsongkey}%
    {\noexpand\the\noexpand\numexpr\noexpand\c@page-\griffsongstart\noexpand\relax}%
    {#1}%
    {\noexpand\the\noexpand\numexpr\noexpand\pdflastypos+#2\noexpand\relax}%
    {\noexpand\the\noexpand\numexpr\noexpand\pdflastypos-#3\noexpand\relax}}}%
  \grifftmpwrite
}
\newcommand{\griffzeilenmarke}{%
  \ifgriffeaus\else % Griffe ausgeblendet -> keine Positionen schreiben
  \ifhmode
    \griffschreiben
      {\noexpand\the\noexpand\numexpr\noexpand\pdflastxpos
        -\noexpand\ifodd\noexpand\c@page\space\griffleftodd
         \noexpand\else\griffleftevn\noexpand\fi\noexpand\relax}%
      {\griffzeileoben}{\griffzeileunten}%
  \fi
  \fi
}
% Zeilenende in alltt (aktives ^^M) um die Markierung erweitern
\begingroup\catcode`\^^M=\active%
\gdef\griffzeilenhook{\let\griffaltcr^^M\def^^M{\griffzeilenmarke\griffaltcr}\let\griffcrtoken^^M}%
\endgroup

% --- Bausteine ---
\newcommand{\griffgrau}{% nur in \griffe grau, Grifftabelle bleibt unverändert
  \colorlet{inkheading}{inkmeta}\colorlet{inkmain}{inkmeta}%
  \colorlet{inkchords}{inkmeta}\colorlet{rulecolor}{inkmeta}%
}
\newcommand{\griffeinzelnbild}[1]{%
  \ifcsdef{griffdef:\detokenize{#1}}%
    {\csuse{griffdef:\detokenize{#1}}}%
    {\PackageWarning{style}{Griff '\detokenize{#1}' nicht definiert}%
     \griffzeichnen{#1}{.,.,.,.,.,.}{}}%
}

% --- Benutzerbefehl ---
\newcommand{\griffe}{%
  \begingroup
  % alltt: Leerzeichen/Zeilenenden und aktive Zeichen normalisieren
  \catcode`\ =10 \catcode`\^^M=10
  \catcode`\,=12 \catcode`\-=12 \catcode`\'=12 \catcode`\`=12
  \catcode`\<=12 \catcode`\>=12
  \@ifnextchar[{\griffeintern}{\griffeintern[0]}}
\def\griffeintern[#1]#2{%
  \endgroup
  \let\griff@danach\relax
  \ifgriffeaus
    % ausgeblendet: Befehl verschwindet spurlos (steht er in einer eigenen
    % Zeile, wird diese geschluckt -> keine Leerzeile)
    \ifvmode\def\griff@danach{\futurelet\griff@next\griff@crweg}\fi
  \else
    % eigene Zeile -> wird (wie bei \schlagmuster) zur Leerzeile
    \leavevmode\griffesetzen{#1}{#2}%
  \fi
  \griff@danach
}
% Zeilenende nach einem Befehl schlucken (wird von \anmerkung benutzt)
\def\griff@crweg{\ifx\griff@next\griffcrtoken\expandafter\@gobble\fi}

% am rechten Textrand auf Höhe der aktuellen Zeile überlagern (wie \schlagmuster)
\newsavebox{\griffueberlagerbox}
\newcommand{\griffeueberlagern}{%
  % Fläche als Hindernis merken (z.B. für eine \anmerkung daneben)
  \griffschreiben{1073741823}%
    {\number\dimexpr0.8\baselineskip\relax}%
    {\number\dimexpr\ht\griffueberlagerbox+\dp\griffueberlagerbox-0.8\baselineskip\relax}%
  \begin{tikzpicture}[remember picture,overlay]
    \node[anchor=north east,inner sep=0pt]
      at ($(current page text area.east |- 0,0)+(0,0.8\baselineskip)$)
      {\usebox{\griffueberlagerbox}};
  \end{tikzpicture}%
}
\makeatother

\ExplSyntaxOn
\cs_set_protected:Npn \griffhindernis #1#2#3#4#5
  {
    \seq_if_exist:cF { g__griff_#1_#2_seq } { \seq_new:c { g__griff_#1_#2_seq } }
    \seq_gput_right:cn { g__griff_#1_#2_seq } { {#3}{#4}{#5} }
  }
\int_new:N \l__griff_n_int
\int_new:N \l__griff_k_int
\int_new:N \l__griff_rows_int
\int_new:N \l__griff_limit_int
\int_new:N \l__griff_h_int
\int_new:N \l__griff_ymin_int
\int_new:N \l__griff_ymax_int
\int_new:N \l__griff_best_int
\int_new:N \l__griff_bestk_int
\int_new:N \l__griff_cbest_int
\int_new:N \l__griff_vgap_int
\dim_new:N \l__griff_w_dim
\dim_new:N \g__griff_x_dim
\dim_new:N \g__griff_y_dim
\box_new:N \l__griff_row_box
\box_new:N \g__griff_block_box
\box_new:N \g__griff_flow_box
\box_new:N \g__griff_place_box
\bool_new:N \l__griff_ok_bool
\tl_new:N \l__griff_seq_tl
\tl_new:N \l__griff_scale_tl
\tl_set:Nn \l__griff_scale_tl { 1 }
\msg_new:nnn { griffe } { kein-platz }
  { Griffe~in~Song~#1~passen~nicht~mehr~unter~den~Text~-~besser~den~Befehl~
    hinten~an~eine~Zeile~hängen~(dann~rechts~am~Rand). }

% Einzelne Griffbilder in Boxen g__griff_item_<i>_box
\cs_new_protected:Npn \__griff_items:n #1
  {
    \int_zero:N \l__griff_n_int
    \clist_map_inline:nn {#1}
      {
        \int_incr:N \l__griff_n_int
        \box_if_exist:cF { g__griff_item_ \int_use:N \l__griff_n_int _box }
          { \box_new:c { g__griff_item_ \int_use:N \l__griff_n_int _box } }
        \hbox_gset:cn { g__griff_item_ \int_use:N \l__griff_n_int _box }
          { \griffgrau \griffeinzelnbild {##1} }
      }
  }
% Zeile #1 (1-basiert) bei #2 Griffen pro Zeile -> \l__griff_row_box
\cs_new_protected:Npn \__griff_row:nn #1#2
  {
    \hbox_set:Nn \l__griff_row_box
      {
        \int_step_inline:nnn { (#1 - 1) * #2 + 1 }
          { \int_min:nn { #1 * #2 } { \l__griff_n_int } }
          {
            \int_compare:nNnT {##1} > { (#1 - 1) * #2 + 1 }
              { \skip_horizontal:n { \griffeabstand } }
            \box_use:c { g__griff_item_ ##1 _box }
          }
      }
  }
% Block mit #1 Griffen pro Zeile (rechtsbündig) -> \g__griff_block_box, Breite -> \l__griff_w_dim
\cs_new_protected:Npn \__griff_build:n #1
  {
    \int_set:Nn \l__griff_rows_int { \int_div_truncate:nn { \l__griff_n_int + #1 - 1 } {#1} }
    \dim_zero:N \l__griff_w_dim
    \int_step_inline:nn { \l__griff_rows_int }
      {
        \__griff_row:nn {##1} {#1}
        \dim_set:Nn \l__griff_w_dim
          { \dim_max:nn { \l__griff_w_dim } { \box_wd:N \l__griff_row_box } }
      }
    \vbox_gset:Nn \g__griff_block_box
      {
        \cs_set_eq:NN \par \tex_par:D % alltt-\par würde eine Leerzeile erzeugen
        \int_step_inline:nn { \l__griff_rows_int }
          {
            \int_compare:nNnT {##1} > {1} { \skip_vertical:n { \griffeabstand } }
            \__griff_row:nn {##1} {#1}
            \hbox_to_wd:nn { \l__griff_w_dim } { \tex_hfil:D \box_use:N \l__griff_row_box }
          }
      }
    \tl_if_eq:NnF \l__griff_scale_tl { 1 }
      {
        \box_gscale:Nnn \g__griff_block_box { \l__griff_scale_tl } { \l__griff_scale_tl }
        \dim_set:Nn \l__griff_w_dim { \box_wd:N \g__griff_block_box }
      }
  }

% --- unter dem Text ---
\cs_new_protected:Npn \__griff_flow:
  {
    \mode_if_horizontal:T { \tex_unskip:D \tex_par:D }
    \skip_vertical:n { \griffeoben }
    \tex_noindent:D \box_use:N \g__griff_flow_box \tex_par:D
  }

% --- neben dem Text ---
\cs_new_protected:Npn \__griff_ymax:nnn #1#2#3
  { \int_set:Nn \l__griff_ymax_int { \int_max:nn { \l__griff_ymax_int } {#2} } }
\cs_new_protected:Npn \__griff_check:nnnn #1#2#3#4 % Kandidat, x, oben, unten
  {
    \bool_if:nT
      {
        \int_compare_p:nNn {#2} > { \l__griff_limit_int } &&
        \int_compare_p:nNn { #3 + \l__griff_vgap_int } > {#1} &&
        \int_compare_p:nNn { #4 - \l__griff_vgap_int } < { #1 + \l__griff_h_int }
      }
      { \bool_set_false:N \l__griff_ok_bool \seq_map_break: }
  }
\cs_new_protected:Npn \__griff_cand:n #1 % Unterkante-Kandidat
  {
    \bool_if:nT
      {
        \int_compare_p:nNn {#1} < { \l__griff_cbest_int } &&
        ! \int_compare_p:nNn {#1} < { \l__griff_ymin_int } &&
        ! \int_compare_p:nNn { #1 + \l__griff_h_int } > { \l__griff_ymax_int }
      }
      {
        \bool_set_true:N \l__griff_ok_bool
        \seq_map_inline:cn { \l__griff_seq_tl } { \__griff_check:nnnn {#1} ##1 }
        \bool_if:NT \l__griff_ok_bool { \int_set:Nn \l__griff_cbest_int {#1} }
      }
  }
\cs_new_protected:Npn \__griff_cand_from:nnn #1#2#3
  {
    \int_compare:nNnT {#1} > { \l__griff_limit_int }
      { \__griff_cand:n { \int_eval:n { #2 + \l__griff_vgap_int } } }
  }
% Suche vorbereiten: Positionsdaten der aktuellen Seite, untere/obere Grenze.
% Ergebnis in \l__griff_setup_bool (false = keine Daten, z.B. erster Lauf)
\bool_new:N \l__griff_setup_bool
\cs_new_protected:Npn \__griff_setup:
  {
    \bool_set_false:N \l__griff_setup_bool
    \tl_set:Nx \l__griff_seq_tl
      { g__griff_ \griffsongkey _ \int_eval:n { \value{page} - \griffsongstart } _seq }
    \seq_if_exist:cT { \l__griff_seq_tl }
      {
        \bool_set_true:N \l__griff_setup_bool
        \int_set:Nn \l__griff_vgap_int { \dim_to_decimal_in_sp:n { 1.5mm } }
        % Seitenzahl unten freihalten (immer, damit das Ergebnis nicht von
        % gerader/ungerader Seite abhängt -> stabile Seitenzahl je Song)
        % ... bzw. Platz für eine Anmerkung (\anmerkung) unten auf der Seite
        \int_set:Nn \l__griff_ymin_int
          {
            \griffbottom + \int_max:nn { 786432 }
              { \dim_to_decimal_in_sp:n { \l__note_need_dim } }
          }
        \int_set:Nn \l__griff_ymax_int { - \c_max_int }
        \seq_map_inline:cn { \l__griff_seq_tl } { \__griff_ymax:nnn ##1 }
      }
  }
% tiefste freie Unterkante für einen Block (Breite/Höhe in
% \l__griff_limit_int / \l__griff_h_int) -> \l__griff_cbest_int
\cs_new_protected:Npn \__griff_find:
  {
    \int_set:Nn \l__griff_cbest_int { \c_max_int }
    \__griff_cand:n { \int_use:N \l__griff_ymin_int }
    \seq_map_inline:cn { \l__griff_seq_tl } { \__griff_cand_from:nnn ##1 }
  }
% Block bauen: #1 Griffe pro Zeile (0 = automatisch: so viele wie in die
% Breite passen) -> \g__griff_block_box
\cs_new_protected:Npn \__griff_block:nn #1#2
  {
    \__griff_items:n {#2}
    \tl_set:Nn \l__griff_scale_tl { 1 }
    \int_compare:nNnTF {#1} > { 0 }
      { \int_set:Nn \l__griff_k_int { \int_min:nn {#1} { \l__griff_n_int } } }
      { \int_set_eq:NN \l__griff_k_int \l__griff_n_int }
    \__griff_build:n { \l__griff_k_int }
    \int_compare:nNnT {#1} = { 0 }
      {
        \bool_while_do:nn
          {
            \dim_compare_p:nNn { \l__griff_w_dim } > { \linewidth } &&
            \int_compare_p:nNn { \l__griff_k_int } > { 1 }
          }
          { \int_decr:N \l__griff_k_int \__griff_build:n { \l__griff_k_int } }
      }
  }

% Griffe am rechten Rand ÜBERLAGERN (verschiebt den Text nie)
% #1 Griffe pro Zeile (0 = automatisch), #2 Liste
\cs_new_protected:Npn \griffesetzen #1#2
  {
    \__griff_block:nn {#1} {#2}
    \box_gset_eq:NN \griffueberlagerbox \g__griff_block_box
    \griffeueberlagern
  }
